% Part: lambda-calculus
% Chapter: lambda-definability

\documentclass[../../../include/open-logic-chapter]{subfiles}

\begin{document}

\olchapter{lam}{ldf}{લૅમ્બડા દ્વારા વ્યાખ્યાયિતતા}
\sourcecorrection{OLLAM-045}{સ્થિર અંગ્રેજી મૂળના
અધ્યાય-ચાલકમાં ઓળખ~\texttt{rep} છે, પરંતુ આ
અધ્યાયના મોટા ભાગના વિભાગોની ઓળખ~\texttt{ldf}
છે. આંતરિક સંદર્ભો એકરૂપ રહે તે માટે અહીં
અધ્યાયની ઓળખ~\texttt{ldf} લીધી છે; જુદી
ઓળખવાળા વિભાગો તેમના અનુવાદ સમયે અલગ
સુધારવાના છે.}

\begin{editorial}
  આ અધ્યાય પ્રાયોગિક છે. તેમાં વધુ સમજૂતી અને
  ઉદાહરણો જોઈએ. સામગ્રીને વ્યાખ્યાઓ અને
  સાબિતીવાળાં પ્રસ્તાવોમાં વધુ સારી રીતે ગોઠવવી
  જોઈએ.
\end{editorial}

\olimport{introduction}
\olimport{arithmetical-functions}
\olimport{pairs}
\olimport{truth-values}
%\olimport{lists}
\olimport{primitive-recursive-functions}
\olimport{fixpoints}
\olimport{minimization}
\olimport{partial-recursive-functions}
\olimport{lambda-definable-recursive}

\OLEndChapterHook

\end{document}
